\caption{Synthetic tests of plane recovery, reduced gain capture, and the corrected diagnostics. (a) Median largest principal angle $\theta_*=\theta_{\max}(Q,Q_*)$ between the selected and independently planted planes for geometry-only, spectrum-only, plane-rotation, and combined interventions. (b) Median one-step gain-capture fraction $c_Q=\sigma_{\max}(Q^{\mathsf T}FQ)/\sigma_{\max}(F)$ evaluated on the known full operator. Open circles and triangles correspond to refitted M2 and M3; filled grey squares denote the M2 plane held fixed at its baseline orientation. Vertical whiskers in (a,b) show 95\% subject-bootstrap intervals (20 independent synthetic subjects per intervention, pooled across two planted families). The ordinate in (b) is restricted to $[0.62,1.035]$ to resolve differences near full gain capture. (c) Ratios on 30 independently specified, spectrum-matched normal systems, comparing the known normal operators with their finite-data M2-fitted counterparts. The primary version shows $R_A$; the companion version additionally shows $R_\sigma$ in light grey. Small symbols denote individual fitted systems, and thick horizontal segments show medians. The known normal operators yield ratios numerically equal to zero, whereas the fitted ratios are slightly positive but all remain below unity in this control set. (d) Per-subject log contrasts under a spectrum-only planted intervention: $\Delta\log K$, $\Delta\log K_{c,A}$, and $\Delta\log R_A$. Small grey points are the 20 subjects; vertical bold segments indicate the interquartile range and horizontal black ticks indicate medians. The relation $\Delta\log R_A=\Delta\log K-\Delta\log K_{c,A}$ holds per subject; it should not be interpreted as an exact relation among the separate sample medians. Panels (a,b,d) use the T05 ablation benchmark; (c) uses the independent T04 normal-control benchmark. These are controlled synthetic systems and do not establish the behavior of physical non-normality under temporally correlated forcing.}
